%% kapa: exam=1 round=28 year=2017 period=1 form=A subject=economics subject_ko=경제학원론 seq=24 number=64
%% source: Q-Net 감정평가사 기출문제 (cst00309) artlSeq=5199603 file=extracted/1008/5199603/2017년 제28회 감정평가사 1차시험 문제지/1교시 A형.pdf
%% conversion: kordoc 4.24.51 -> tools/gen1.py (auto); review=proofread-sonnet; flags=none
\begin{question}{64}
기술진보가 없는 솔로우(Solow)의 경제성장모형에서 1인당 생산함수는 $y=k^{0.2}$, 저축률은 0.4, 자본의 감가상각률은 0.15, 인구증가율은 0.05이다. 현재 경제가 균제상태(steady state)일 때 다음 중 옳은 것을 모두 고른 것은? (단, $y$는 1인당 생산량, $k$는 1인당 자본량이다.)
\begin{framedtext}
ㄱ. 현재 균제상태의 1인당 자본량은 황금률 수준(golden rule level)의 1인당 자본량보다 작다. \par
ㄴ. 황금률을 달성시키는 저축률은 0.2이다. \par
ㄷ. 인구증가율이 증가하면 황금률 수준의 1인당 자본량도 증가한다. \par
ㄹ. 감가상각률이 증가하면 황금률 수준의 1인당 자본량은 감소한다.
\end{framedtext}
\begin{choices}
  \item ㄱ, ㄴ
  \item ㄱ, ㄷ
  \item ㄴ, ㄹ
  \item ㄱ, ㄴ, ㄹ
  \item ㄴ, ㄷ, ㄹ
\end{choices}
\end{question}
\answer{3}
